\subsection{Substitution and evaluation of power series}

\subsubsection{Defining substitution}

Definition \ref{def.pol.subs} shows that if $f\in K\left[  x\right]  $ is a
polynomial, then almost anything (to be more precise: any element of a
$K$-algebra) can be substituted into $f$.

In contrast, if $f\in K\left[  \left[  x\right]  \right]  $ is a FPS, then
there are far fewer things that can be substituted into $f$. Even elements of
$K$ itself cannot always be substituted into $f$. For example, if we try to
substitute $1$ for $x$ in the FPS $1+x+x^{2}+x^{3}+\cdots$, then we get%
\[
1+1+1^{2}+1^{3}+\cdots=1+1+1+1+\cdots,
\]
which is undefined. Real analysis can help make sense of certain values of
FPSs (for example, substituting $\dfrac{1}{2}$ for $x$ into the FPS
$1+x+x^{2}+x^{3}+\cdots$ yields the convergent series $1+\dfrac{1}{2}%
+\dfrac{1}{2^{2}}+\dfrac{1}{2^{3}}+\cdots=2$), but this is subtle and specific
to certain numbers and certain FPSs.\footnote{For instance, it is not hard to
see that there is no nonzero complex number that can be substituted into the
FPS $\sum_{n\in\mathbb{N}}n!x^{n}$ to obtain a convergent result. Thus, even
though some complex numbers can be substituted into some FPSs, there is no
complex number other than $0$ that can be substituted into \textbf{every}
FPS.}

Thus, polynomials have a clear advantage over FPSs.

However, let us not give up on FPSs yet. \textbf{Some} things can be
substituted into an FPS. For example:

\begin{itemize}
\item We can always substitute $0$ for $x$ in an FPS $a_{0}+a_{1}x+a_{2}%
x^{2}+a_{3}x^{3}+\cdots$. The result is%
\[
a_{0}+a_{1}0+a_{2}0^{2}+a_{3}0^{3}+\cdots=a_{0}+0+0+0+\cdots=a_{0}.
\]


\item We can always substitute $x$ for $x$ in an FPS $a_{0}+a_{1}x+a_{2}%
x^{2}+a_{3}x^{3}+\cdots$. The result is the same FPS $a_{0}+a_{1}x+a_{2}%
x^{2}+a_{3}x^{3}+\cdots$ that we started with (obviously).

\item We can always substitute $2x$ for $x$ in an FPS $a_{0}+a_{1}x+a_{2}%
x^{2}+a_{3}x^{3}+\cdots$. The result is%
\[
a_{0}+a_{1}\left(  2x\right)  +a_{2}\left(  2x\right)  ^{2}+a_{3}\left(
2x\right)  ^{3}+\cdots=a_{0}+2a_{1}x+4a_{2}x^{2}+8a_{3}x^{3}+\cdots.
\]


\item We can always substitute $x^{2}+x$ for $x$ in an FPS $a_{0}+a_{1}%
x+a_{2}x^{2}+a_{3}x^{3}+\cdots$. This is less obvious, so let me explain why.
If we try to substitute $x^{2}+x$ for $x$ in an FPS $a_{0}+a_{1}x+a_{2}%
x^{2}+a_{3}x^{3}+\cdots$, then we obtain%
\begin{align*}
&  \left(  a_{0}+a_{1}x+a_{2}x^{2}+a_{3}x^{3}+\cdots\right)  \left[
x+x^{2}\right] \\
&  =a_{0}+a_{1}\left(  x+x^{2}\right)  +a_{2}\left(  x+x^{2}\right)
^{2}+a_{3}\left(  x+x^{2}\right)  ^{3}+\cdots\\
&  =a_{0}+a_{1}\left(  x+x^{2}\right)  +a_{2}\left(  x^{2}+2x^{3}%
+x^{4}\right)  +a_{3}\left(  x^{3}+3x^{4}+3x^{5}+x^{6}\right)  +\cdots\\
&  =a_{0}+a_{1}x+\left(  a_{1}+a_{2}\right)  x^{2}+\left(  2a_{2}%
+a_{3}\right)  x^{3}+\left(  a_{2}+3a_{3}+a_{4}\right)  x^{4}+\cdots.
\end{align*}
I claim that the right hand side here is well-defined. To prove this, I need
to show that for each $n\in\mathbb{N}$, the coefficient of $x^{n}$ on this
right hand side is a \textbf{finite} sum of $a_{i}$'s. Indeed, fix
$n\in\mathbb{N}$. Recall that the right hand side is obtained by expanding the
infinite sum%
\[
a_{0}+a_{1}\left(  x+x^{2}\right)  +a_{2}\left(  x+x^{2}\right)  ^{2}%
+a_{3}\left(  x+x^{2}\right)  ^{3}+\cdots.
\]
Only the first $n+1$ addends of this infinite sum (i.e., only the addends
$a_{k}\left(  x+x^{2}\right)  ^{k}$ with $k\leq n$) can contribute to the
coefficient of $x^{n}$, since any of the remaining addends is a multiple of
$x^{n+1}$ (because it has the form $a_{k}\left(  x+x^{2}\right)  ^{k}%
=a_{k}\left(  x\left(  1+x\right)  \right)  ^{k}=a_{k}x^{k}\left(  1+x\right)
^{k}$ with $k\geq n+1$) and thus has a zero coefficient of $x^{n}$. Hence, the
coefficient of $x^{n}$ in this infinite sum equals the coefficient of $x^{n}$
in the \textbf{finite} sum
\[
a_{0}+a_{1}\left(  x+x^{2}\right)  +a_{2}\left(  x+x^{2}\right)  ^{2}%
+a_{3}\left(  x+x^{2}\right)  ^{3}+\cdots+a_{n}\left(  x+x^{2}\right)  ^{n}.
\]
But the latter coefficient is clearly a \textbf{finite} sum of $a_{i}$'s.
Thus, my claim is proved, and it follows that the result of substituting
$x^{2}+x$ for $x$ in an FPS $a_{0}+a_{1}x+a_{2}x^{2}+a_{3}x^{3}+\cdots$ is well-defined.
\end{itemize}

The idea of the last example can be generalized; there was nothing special
about $x+x^{2}$ that we used other than the fact that $x+x^{2}$ is a multiple
of $x$ (that is, an FPS whose constant term is $0$). Thus, generalizing our
reasoning from this example, we can convince ourselves that any FPS $g$ that
is a multiple of $x$ (that is, whose constant term is $0$) can be substituted
into any FPS. Let us introduce a notation for this, exactly like we did for
substituting things into polynomials:

\begin{definition}
\label{def.fps.subs}Let $f$ and $g$ be two FPSs in $K\left[  \left[  x\right]
\right]  $. Assume that $\left[  x^{0}\right]  g=0$ (that is, $g=g_{1}%
x^{1}+g_{2}x^{2}+g_{3}x^{3}+\cdots$ for some $g_{1},g_{2},g_{3},\ldots\in K$).

We then define an FPS $f\left[  g\right]  \in K\left[  \left[  x\right]
\right]  $ as follows:

Write $f$ in the form $f=\sum_{n\in\mathbb{N}}f_{n}x^{n}$ with $f_{0}%
,f_{1},f_{2},\ldots\in K$. (That is, $f_{n}=\left[  x^{n}\right]  f$ for each
$n\in\mathbb{N}$.) Then, set%
\begin{equation}
f\left[  g\right]  :=\sum_{n\in\mathbb{N}}f_{n}g^{n}.
\label{eq.def.fps.subs.eq}%
\end{equation}
(This sum is well-defined, as we will see in Proposition
\ref{prop.fps.subs.wd} \textbf{(b)} below.)

This FPS $f\left[  g\right]  $ is also denoted by $f\circ g$, and is called
the \emph{composition} of $f$ with $g$, or the result of \emph{substituting}
$g$ for $x$ in $f$.
\end{definition}

Once again, it is not uncommon to see this FPS $f\left[  g\right]  $ denoted
by $f\left(  g\right)  $, but I will eschew the latter notation (since it can
be confused with a product).

In order to prove that Definition \ref{def.fps.subs} makes sense, we need to
ensure that the infinite sum $\sum_{n\in\mathbb{N}}f_{n}g^{n}$ in
(\ref{eq.def.fps.subs.eq}) is well-defined. The proof of this fact is
analogous to the reasoning I used in the last example; let me present it again
in the general case:

\begin{proposition}
\label{prop.fps.subs.wd}Let $f$ and $g$ be two FPSs in $K\left[  \left[
x\right]  \right]  $. Assume that $\left[  x^{0}\right]  g=0$. Write $f$ in
the form $f=\sum_{n\in\mathbb{N}}f_{n}x^{n}$ with $f_{0},f_{1},f_{2},\ldots\in
K$. Then: \medskip

\textbf{(a)} For each $n\in\mathbb{N}$, the first $n$ coefficients of the FPS
$g^{n}$ are $0$. \medskip

\textbf{(b)} The sum $\sum_{n\in\mathbb{N}}f_{n}g^{n}$ is well-defined, i.e.,
the family $\left(  f_{n}g^{n}\right)  _{n\in\mathbb{N}}$ is summable.
\medskip

\textbf{(c)} We have $\left[  x^{0}\right]  \left(  \sum_{n\in\mathbb{N}}%
f_{n}g^{n}\right)  =f_{0}$.
\end{proposition}

\begin{proof}
[Proof of Proposition \ref{prop.fps.subs.wd}.]\textbf{(a)} This is easily
proved by induction on $n$. Here is a shorter alternative argument:

The FPS $g$ has constant term $\left[  x^{0}\right]  g=0$. Hence, Lemma
\ref{lem.fps.g=xh} (applied to $a=g$) yields that there exists an $h\in
K\left[  \left[  x\right]  \right]  $ such that $g=xh$. Consider this $h$.

Now, let $n\in\mathbb{N}$. From $g=xh$, we obtain $g^{n}=\left(  xh\right)
^{n}=x^{n}h^{n}$. However, Lemma \ref{lem.fps.first-n-coeffs-of-xna} (applied
to $k=n$ and $a=h^{n}$) yields that the first $n$ coefficients of the FPS
$x^{n}h^{n}$ are $0$. In other words, the first $n$ coefficients of the FPS
$g^{n}$ are $0$ (since $g^{n}=x^{n}h^{n}$). Thus, Proposition
\ref{prop.fps.subs.wd} \textbf{(a)} is proved. \medskip

\textbf{(b)} This follows from part \textbf{(a)}. Here are the details.

We must prove that the family $\left(  f_{n}g^{n}\right)  _{n\in\mathbb{N}}$
is summable. In other words, we must prove that the family $\left(  f_{i}%
g^{i}\right)  _{i\in\mathbb{N}}$ is summable (since $\left(  f_{i}%
g^{i}\right)  _{i\in\mathbb{N}}=\left(  f_{n}g^{n}\right)  _{n\in\mathbb{N}}%
$). In other words, we must prove that for each $n\in\mathbb{N}$, all but
finitely many $i\in\mathbb{N}$ satisfy $\left[  x^{n}\right]  \left(
f_{i}g^{i}\right)  =0$ (by the definition of \textquotedblleft
summable\textquotedblright). So let us prove this.

Fix $n\in\mathbb{N}$. We must prove that all but finitely many $i\in
\mathbb{N}$ satisfy $\left[  x^{n}\right]  \left(  f_{i}g^{i}\right)  =0$.

Indeed, let $i\in\mathbb{N}$ satisfy $i>n$. Then, $n<i$. Now, the first $i$
coefficients of the FPS $g^{i}$ are $0$ (by Proposition \ref{prop.fps.subs.wd}
\textbf{(a)}, applied to $i$ instead of $n$). However, the coefficient
$\left[  x^{n}\right]  \left(  g^{i}\right)  $ of $g^{i}$ is one of these
first $i$ coefficients (because $n<i$). Thus, this coefficient $\left[
x^{n}\right]  \left(  g^{i}\right)  $ must be $0$. Now, $f_{i}\in K$; thus,
(\ref{pf.thm.fps.ring.xn(la)=}) (applied to $\lambda=f_{i}$ and $\mathbf{a}%
=g^{i}$) yields $\left[  x^{n}\right]  \left(  f_{i}g^{i}\right)  =f_{i}%
\cdot\underbrace{\left[  x^{n}\right]  \left(  g^{i}\right)  }_{=0}=0$.

Forget that we fixed $i$. We thus have shown that all $i\in\mathbb{N}$
satisfying $i>n$ satisfy $\left[  x^{n}\right]  \left(  f_{i}g^{i}\right)
=0$. Hence, all but finitely many $i\in\mathbb{N}$ satisfy $\left[
x^{n}\right]  \left(  f_{i}g^{i}\right)  =0$ (because all but finitely many
$i\in\mathbb{N}$ satisfy $i>n$). This is precisely what we wanted to prove.
Thus, Proposition \ref{prop.fps.subs.wd} \textbf{(b)} is proved. \medskip

\textbf{(c)} Let $n$ be a positive integer. We shall first show that $\left[
x^{0}\right]  \left(  f_{n}g^{n}\right)  =0$.

Indeed, Proposition \ref{prop.fps.subs.wd} \textbf{(a)} shows that the first
$n$ coefficients of the FPS $g^{n}$ are $0$. However, the coefficient $\left[
x^{0}\right]  \left(  g^{n}\right)  $ is one of these first $n$ coefficients
(since $n$ is positive). Thus, this coefficient $\left[  x^{0}\right]  \left(
g^{n}\right)  $ must be $0$. Now, $f_{n}\in K$; thus,
(\ref{pf.thm.fps.ring.xn(la)=}) (applied to $f_{n}$, $g^{n}$ and $0$ instead
of $\lambda$, $\mathbf{a}$ and $n$) yields $\left[  x^{0}\right]  \left(
f_{n}g^{n}\right)  =f_{n}\cdot\underbrace{\left[  x^{0}\right]  \left(
g^{n}\right)  }_{=0}=0$.

Forget that we fixed $n$. We thus have shown that
\begin{equation}
\left[  x^{0}\right]  \left(  f_{n}g^{n}\right)
=0\ \ \ \ \ \ \ \ \ \ \text{for each positive integer }n.
\label{pf.prop.fps.subs.wd.c.1}%
\end{equation}


Now,%
\begin{align*}
\left[  x^{0}\right]  \left(  \sum_{n\in\mathbb{N}}f_{n}g^{n}\right)   &
=\sum_{n\in\mathbb{N}}\left[  x^{0}\right]  \left(  f_{n}g^{n}\right)
\ \ \ \ \ \ \ \ \ \ \left(  \text{by (\ref{eq.def.fps.summable.sum})}\right)
\\
&  =\left[  x^{0}\right]  \left(  f_{0}\underbrace{g^{0}}_{=\underline{1}%
}\right)  +\sum_{n>0}\underbrace{\left[  x^{0}\right]  \left(  f_{n}%
g^{n}\right)  }_{\substack{=0\\\text{(by (\ref{pf.prop.fps.subs.wd.c.1}))}}}\\
&  \ \ \ \ \ \ \ \ \ \ \ \ \ \ \ \ \ \ \ \ \left(
\begin{array}
[c]{c}%
\text{here, we have split off the}\\
\text{addend for }n=0\text{ from the sum}%
\end{array}
\right) \\
&  =\left[  x^{0}\right]  \underbrace{\left(  f_{0}\underline{1}\right)
}_{\substack{=f_{0}\left(  1,0,0,0,\ldots\right)  \\=\left(  f_{0}%
,0,0,0,\ldots\right)  }}+\underbrace{\sum_{n>0}0}_{=0}=\left[  x^{0}\right]
\left(  f_{0},0,0,0,\ldots\right)  =f_{0}.
\end{align*}
This proves Proposition \ref{prop.fps.subs.wd} \textbf{(c)}.
\end{proof}

\begin{example}
\label{exa.fps.subs.fibonacci}The FPS $x+x^{2}$ has constant term $\left[
x^{0}\right]  \left(  x+x^{2}\right)  =0$. Hence, according to Definition
\ref{def.fps.subs}, we can substitute it for $x$ into $1+x+x^{2}+x^{3}+\cdots
$. The result is%
\begin{align*}
&  \left(  1+x+x^{2}+x^{3}+\cdots\right)  \left[  x+x^{2}\right] \\
&  =1+\left(  x+x^{2}\right)  +\left(  x+x^{2}\right)  ^{2}+\left(
x+x^{2}\right)  ^{3}+\left(  x+x^{2}\right)  ^{4}+\left(  x+x^{2}\right)
^{5}+\cdots\\
&  =1+x+2x^{2}+3x^{3}+5x^{4}+8x^{5}+\cdots.
\end{align*}
The right hand side appears to be $f_{1}+f_{2}x+f_{3}x^{2}+f_{4}x^{3}+\cdots$,
where $\left(  f_{0},f_{1},f_{2},\ldots\right)  $ is the Fibonacci sequence
(as defined in Section \ref{sec.gf.exas}). Let me show that this indeed the case.

In Example 1 in Section \ref{sec.gf.exas}, we had shown that%
\[
f_{0}+f_{1}x+f_{2}x^{2}+f_{3}x^{3}+\cdots=\dfrac{x}{1-x-x^{2}}.
\]
Thus,%
\begin{align*}
\dfrac{x}{1-x-x^{2}}  &  =\underbrace{f_{0}}_{=0}+\,f_{1}x+f_{2}x^{2}%
+f_{3}x^{3}+\cdots\\
&  =f_{1}x+f_{2}x^{2}+f_{3}x^{3}+f_{4}x^{4}+\cdots\\
&  =x\left(  f_{1}+f_{2}x+f_{3}x^{2}+f_{4}x^{3}+\cdots\right)  .
\end{align*}
Cancelling $x$ from this equality (this is indeed allowed -- make sure you
understand why!), we obtain%
\[
\dfrac{1}{1-x-x^{2}}=f_{1}+f_{2}x+f_{3}x^{2}+f_{4}x^{3}+\cdots.
\]
However, it appears reasonable to expect that
\begin{equation}
\dfrac{1}{1-x-x^{2}}=\dfrac{1}{1-x}\left[  x+x^{2}\right]  ,
\label{eq.exa.fps.subs.fibonacci.subs-plausible}%
\end{equation}
because substituting $x+x^{2}$ for $x$ in the expression $\dfrac{1}{1-x}$
results in $\dfrac{1}{1-x-x^{2}}$. This is plausible but not obvious -- after
all, we defined $\dfrac{1}{1-x}\left[  x+x^{2}\right]  $ to be the result of
substituting $x+x^{2}$ for $x$ into the \textbf{expanded} version of
$\dfrac{1}{1-x}$ (which is $1+x+x^{2}+x^{3}+\cdots$), not into the fractional
expression $\dfrac{1}{1-x}$. Nevertheless,
(\ref{eq.exa.fps.subs.fibonacci.subs-plausible}) is true (and will soon be
proved). If we take this fact for granted, then our claim easily follows:%
\begin{align*}
f_{1}+f_{2}x+f_{3}x^{2}+f_{4}x^{3}+\cdots &  =\dfrac{1}{1-x-x^{2}}=\dfrac
{1}{1-x}\left[  x+x^{2}\right] \\
&  =\left(  1+x+x^{2}+x^{3}+\cdots\right)  \left[  x+x^{2}\right]
\end{align*}
(since $\dfrac{1}{1-x}=1+x+x^{2}+x^{3}+\cdots$).
\end{example}

\subsubsection{Laws of substitution}

The plausible but nontrivial statement
(\ref{eq.exa.fps.subs.fibonacci.subs-plausible}) that we have just used
follows from part \textbf{(c)} of the following proposition:\footnote{We are
treating the symbol \textquotedblleft$\circ$\textquotedblright\ similarly to
the multiplication sign $\cdot$ in our PEMDAS convention. Thus, an expression
like \textquotedblleft$f_{1}\circ g+f_{2}\circ g$\textquotedblright\ is
understood to mean $\left(  f_{1}\circ g\right)  +\left(  f_{2}\circ g\right)
$.}

\begin{proposition}
\label{prop.fps.subs.rules}Composition of FPSs satisfies the rules you would
expect it to satisfy: \medskip

\textbf{(a)} If $f_{1},f_{2},g\in K\left[  \left[  x\right]  \right]  $
satisfy $\left[  x^{0}\right]  g=0$, then $\left(  f_{1}+f_{2}\right)  \circ
g=f_{1}\circ g+f_{2}\circ g$. \medskip

\textbf{(b)} If $f_{1},f_{2},g\in K\left[  \left[  x\right]  \right]  $
satisfy $\left[  x^{0}\right]  g=0$, then $\left(  f_{1}\cdot f_{2}\right)
\circ g=\left(  f_{1}\circ g\right)  \cdot\left(  f_{2}\circ g\right)  $.
\medskip

\textbf{(c)} If $f_{1},f_{2},g\in K\left[  \left[  x\right]  \right]  $
satisfy $\left[  x^{0}\right]  g=0$, then $\dfrac{f_{1}}{f_{2}}\circ
g=\dfrac{f_{1}\circ g}{f_{2}\circ g}$, as long as $f_{2}$ is invertible. (In
particular, $f_{2}\circ g$ is automatically invertible under these
assumptions.) \medskip

\textbf{(d)} If $f,g\in K\left[  \left[  x\right]  \right]  $ satisfy $\left[
x^{0}\right]  g=0$, then $f^{k}\circ g=\left(  f\circ g\right)  ^{k}$ for each
$k\in\mathbb{N}$. \medskip

\textbf{(e)} If $f,g,h\in K\left[  \left[  x\right]  \right]  $ satisfy
$\left[  x^{0}\right]  g=0$ and $\left[  x^{0}\right]  h=0$, then $\left[
x^{0}\right]  \left(  g\circ h\right)  =0$ and $\left(  f\circ g\right)  \circ
h=f\circ\left(  g\circ h\right)  $. \medskip

\textbf{(f)} We have $\underline{a}\circ g=\underline{a}$ for each $a\in K$
and $g\in K\left[  \left[  x\right]  \right]  $. \medskip

\textbf{(g)} We have $x\circ g=g\circ x=g$ for each $g\in K\left[  \left[
x\right]  \right]  $. \medskip

\textbf{(h)} If $\left(  f_{i}\right)  _{i\in I}\in K\left[  \left[  x\right]
\right]  ^{I}$ is a summable family of FPSs, and if $g\in K\left[  \left[
x\right]  \right]  $ is an FPS satisfying $\left[  x^{0}\right]  g=0$, then
the family $\left(  f_{i}\circ g\right)  _{i\in I}\in K\left[  \left[
x\right]  \right]  ^{I}$ is summable as well and we have $\left(  \sum_{i\in
I}f_{i}\right)  \circ g=\sum_{i\in I}f_{i}\circ g$.
\end{proposition}

For our proof of Proposition \ref{prop.fps.subs.rules}, we will need the
following lemma:

\begin{lemma}
\label{lem.fps.fg-coeffs-0}Let $f,g\in K\left[  \left[  x\right]  \right]  $
satisfy $\left[  x^{0}\right]  g=0$. Let $k\in\mathbb{N}$ be such that the
first $k$ coefficients of $f$ are $0$. Then, the first $k$ coefficients of
$f\circ g$ are $0$.
\end{lemma}

\begin{proof}
[Proof of Lemma \ref{lem.fps.fg-coeffs-0}.]This is very similar to the proof
of Proposition \ref{prop.fps.subs.wd} \textbf{(a)}.

We have $\left[  x^{0}\right]  g=0$. Hence, Lemma \ref{lem.fps.g=xh} (applied
to $a=g$) yields that there exists an $h\in K\left[  \left[  x\right]
\right]  $ such that $g=xh$. Consider this $h$.

Write the FPS $f$ in the form $f=\left(  f_{0},f_{1},f_{2},\ldots\right)  $.
Then, the first $k$ coefficients of $f$ are $f_{0},f_{1},\ldots,f_{k-1}$.
Hence, these coefficients $f_{0},f_{1},\ldots,f_{k-1}$ are $0$ (since the
first $k$ coefficients of $f$ are $0$). In other words,%
\begin{equation}
f_{n}=0\ \ \ \ \ \ \ \ \ \ \text{for each }n<k.
\label{pf.lem.fps.fg-coeffs-0.gn=0}%
\end{equation}


Now, $f=\left(  f_{0},f_{1},f_{2},\ldots\right)  =\sum_{n\in\mathbb{N}}%
f_{n}x^{n}$ with $f_{0},f_{1},f_{2},\ldots\in K$. Hence, Definition
\ref{def.fps.subs} yields%
\begin{align*}
f\left[  g\right]   &  =\sum_{n\in\mathbb{N}}f_{n}g^{n}=\sum_{\substack{n\in
\mathbb{N};\\n<k}}\underbrace{f_{n}}_{\substack{=0\\\text{(by
(\ref{pf.lem.fps.fg-coeffs-0.gn=0}))}}}g^{n}+\underbrace{\sum_{\substack{n\in
\mathbb{N};\\n\geq k}}}_{=\sum_{\substack{n\in\mathbb{N};\\k\leq n}}}%
f_{n}\underbrace{g^{n}}_{\substack{=\left(  xh\right)  ^{n}\\\text{(since
}g=xh\text{)}}}=\underbrace{\sum_{\substack{n\in\mathbb{N};\\n<k}}0g^{n}}%
_{=0}+\sum_{\substack{n\in\mathbb{N};\\k\leq n}}f_{n}\left(  xh\right)  ^{n}\\
&  =\sum_{\substack{n\in\mathbb{N};\\k\leq n}}f_{n}\underbrace{\left(
xh\right)  ^{n}}_{=x^{n}h^{n}}=\sum_{\substack{n\in\mathbb{N};\\k\leq n}%
}f_{n}x^{n}h^{n}.
\end{align*}
Thus,
\[
f\circ g=f\left[  g\right]  =\sum_{\substack{n\in\mathbb{N};\\k\leq
n}}\underbrace{f_{n}x^{n}}_{=x^{n}f_{n}}h^{n}=\sum_{\substack{n\in
\mathbb{N};\\k\leq n}}x^{n}f_{n}h^{n}.
\]
But this ensures that the first $k$ coefficients of $f\circ g$ are
$0$\ \ \ \ \footnote{\textit{Proof.} We must show that $\left[  x^{m}\right]
\left(  f\circ g\right)  =0$ for any nonnegative integer $m<k$. But we can do
this directly: If $m$ is a nonnegative integer such that $m<k$, then%
\begin{align*}
\left[  x^{m}\right]  \left(  f\circ g\right)   &  =\left[  x^{m}\right]
\left(  \sum_{\substack{n\in\mathbb{N};\\k\leq n}}x^{n}f_{n}h^{n}\right)
\ \ \ \ \ \ \ \ \ \ \left(  \text{since }f\circ g=\sum_{\substack{n\in
\mathbb{N};\\k\leq n}}x^{n}f_{n}h^{n}\right) \\
&  =\sum_{\substack{n\in\mathbb{N};\\k\leq n}}\underbrace{\left[
x^{m}\right]  \left(  x^{n}f_{n}h^{n}\right)  }_{\substack{=\sum_{i=0}%
^{m}\left[  x^{i}\right]  \left(  x^{n}\right)  \cdot\left[  x^{m-i}\right]
\left(  f_{n}h^{n}\right)  \\\text{(by (\ref{pf.thm.fps.ring.xn(ab)=2}%
),}\\\text{applied to }x^{n}\text{, }f_{n}h^{n}\text{ and }m\\\text{instead of
}\mathbf{a}\text{, }\mathbf{b}\text{ and }n\text{)}}}=\sum_{\substack{n\in
\mathbb{N};\\k\leq n}}\ \ \sum_{i=0}^{m}\underbrace{\left[  x^{i}\right]
\left(  x^{n}\right)  }_{\substack{=0\\\text{(since }i\leq m<k\leq
n\\\text{and thus }i\neq n\text{)}}}\cdot\left[  x^{m-i}\right]  \left(
f_{n}h^{n}\right) \\
&  =\sum_{\substack{n\in\mathbb{N};\\k\leq n}}\ \ \sum_{i=0}^{m}0\cdot\left[
x^{m-i}\right]  \left(  f_{n}h^{n}\right)  =0,
\end{align*}
exactly as we wanted to show.}. Thus, Lemma \ref{lem.fps.fg-coeffs-0} follows.
\end{proof}

Our proof of Proposition \ref{prop.fps.subs.rules} will furthermore use the
\emph{Kronecker delta notation}:

\begin{definition}
\label{def.kron-delta}If $i$ and $j$ are any objects, then $\delta_{i,j}$
means the element $%
\begin{cases}
1, & \text{if }i=j;\\
0, & \text{if }i\neq j
\end{cases}
$\ \ \ of $K$.
\end{definition}

For example, $\delta_{2,2}=1$ and $\delta_{3,8}=0$.

\begin{proof}
[Proof of Proposition \ref{prop.fps.subs.rules}.]The proof is long and not
particularly combinatorial. I am merely writing it down because it is so
rarely explained in the literature. \medskip

\textbf{(a)} This is an easy consequence of the definitions, and also appears
in \cite[Theorem 7.62]{Loehr-BC} and \cite[Proposition 2.2.2]{Brewer14}.

Here are the details: Let $f_{1},f_{2},g\in K\left[  \left[  x\right]
\right]  $ satisfy $\left[  x^{0}\right]  g=0$. Write the FPSs $f_{1}$ and
$f_{2}$ as
\begin{equation}
f_{1}=\sum_{n\in\mathbb{N}}f_{1,n}x^{n}\ \ \ \ \ \ \ \ \ \ \text{and}%
\ \ \ \ \ \ \ \ \ \ f_{2}=\sum_{n\in\mathbb{N}}f_{2,n}x^{n}
\label{pf.prop.fps.subs.rules.a.1}%
\end{equation}
with $f_{1,0},f_{1,1},f_{1,2},\ldots\in K$ and $f_{2,0},f_{2,1},f_{2,2}%
,\ldots\in K$. Then, adding the two equalities in
(\ref{pf.prop.fps.subs.rules.a.1}) together, we find%
\[
f_{1}+f_{2}=\sum_{n\in\mathbb{N}}f_{1,n}x^{n}+\sum_{n\in\mathbb{N}}%
f_{2,n}x^{n}=\sum_{n\in\mathbb{N}}\underbrace{\left(  f_{1,n}x^{n}%
+f_{2,n}x^{n}\right)  }_{=\left(  f_{1,n}+f_{2,n}\right)  x^{n}}=\sum
_{n\in\mathbb{N}}\left(  f_{1,n}+f_{2,n}\right)  x^{n}.
\]
Thus, Definition \ref{def.fps.subs} (applied to $f=f_{1}+f_{2}$) yields%
\begin{equation}
\left(  f_{1}+f_{2}\right)  \left[  g\right]  =\sum_{n\in\mathbb{N}}\left(
f_{1,n}+f_{2,n}\right)  g^{n} \label{pf.prop.fps.subs.rules.a.3}%
\end{equation}
(since $f_{1,n}+f_{2,n}\in K$ for each $n\in\mathbb{N}$).

On the other hand, we have $f_{1}=\sum_{n\in\mathbb{N}}f_{1,n}x^{n}$. Thus,
Definition \ref{def.fps.subs} yields $f_{1}\left[  g\right]  =\sum
_{n\in\mathbb{N}}f_{1,n}g^{n}$ (since $f_{1,n}\in K$ for each $n\in\mathbb{N}%
$). Similarly, $f_{2}\left[  g\right]  =\sum_{n\in\mathbb{N}}f_{2,n}g^{n}$.
Adding these two equalities together, we obtain%
\[
f_{1}\left[  g\right]  +f_{2}\left[  g\right]  =\sum_{n\in\mathbb{N}}%
f_{1,n}g^{n}+\sum_{n\in\mathbb{N}}f_{2,n}g^{n}=\sum_{n\in\mathbb{N}%
}\underbrace{\left(  f_{1,n}g^{n}+f_{2,n}g^{n}\right)  }_{=\left(
f_{1,n}+f_{2,n}\right)  g^{n}}=\sum_{n\in\mathbb{N}}\left(  f_{1,n}%
+f_{2,n}\right)  g^{n}.
\]
Comparing this with (\ref{pf.prop.fps.subs.rules.a.3}), we obtain $\left(
f_{1}+f_{2}\right)  \left[  g\right]  =f_{1}\left[  g\right]  +f_{2}\left[
g\right]  $. In other words, $\left(  f_{1}+f_{2}\right)  \circ g=f_{1}\circ
g+f_{2}\circ g$ (since the notation $f\circ g$ is synonymous to $f\left[
g\right]  $). This proves Proposition \ref{prop.fps.subs.rules} \textbf{(a)}.

\bigskip

\textbf{(f)} This is a near-trivial consequence of the definitions. To wit:
Let $a\in K$. Then,
\begin{align}
\sum_{n\in\mathbb{N}}a\delta_{n,0}x^{n}  &  =a\underbrace{\delta_{0,0}%
}_{\substack{=1\\\text{(since }0=0\text{)}}}\underbrace{x^{0}}_{=\underline{1}%
}+\sum_{\substack{n\in\mathbb{N};\\n\neq0}}a\underbrace{\delta_{n,0}%
}_{\substack{=0\\\text{(since }n\neq0\text{)}}}x^{n}=a\cdot\underline{1}%
+\underbrace{\sum_{\substack{n\in\mathbb{N};\\n\neq0}}a\cdot0x^{n}}%
_{=0}\nonumber\\
&  =a\cdot\underline{1}=a\cdot\left(  1,0,0,0,\ldots\right)  =\left(
a\cdot1,a\cdot0,a\cdot0,a\cdot0,\ldots\right) \nonumber\\
&  =\left(  a,0,0,0,\ldots\right)  =\underline{a}.
\label{pf.prop.fps.subs.rules.f.1}%
\end{align}


Let $g\in K\left[  \left[  x\right]  \right]  $. We must prove that
$\underline{a}\circ g=\underline{a}$. From (\ref{pf.prop.fps.subs.rules.f.1}),
we obtain $\underline{a}=\sum_{n\in\mathbb{N}}a\delta_{n,0}x^{n}$. Hence,
Definition \ref{def.fps.subs} (or Definition \ref{def.pol.subs}) yields
$\underline{a}\left[  g\right]  =\sum_{n\in\mathbb{N}}a\delta_{n,0}g^{n}$
(since $a\delta_{n,0}\in K$ for each $n\in\mathbb{N}$). Thus,%
\begin{align*}
\underline{a}\left[  g\right]   &  =\sum_{n\in\mathbb{N}}a\delta_{n,0}%
g^{n}=a\underbrace{\delta_{0,0}}_{\substack{=1\\\text{(since }0=0\text{)}%
}}\underbrace{g^{0}}_{=\underline{1}}+\sum_{\substack{n\in\mathbb{N};\\n\neq
0}}a\underbrace{\delta_{n,0}}_{\substack{=0\\\text{(since }n\neq0\text{)}%
}}g^{n}=a\cdot\underline{1}+\underbrace{\sum_{\substack{n\in\mathbb{N}%
;\\n\neq0}}a\cdot0g^{n}}_{=0}\\
&  =a\cdot\underline{1}=\underline{a}.
\end{align*}
In other words, $\underline{a}\circ g=\underline{a}$ (since $\underline{a}%
\circ g$ is a synonym for $\underline{a}\left[  g\right]  $). This proves
Proposition \ref{prop.fps.subs.rules} \textbf{(f)}.

\bigskip

\textbf{(g)} This is easy, too. Indeed, we have%
\begin{align}
\sum_{n\in\mathbb{N}}\delta_{n,1}x^{n}  &  =\underbrace{\delta_{1,1}%
}_{\substack{=1\\\text{(since }1=1\text{)}}}\underbrace{x^{1}}_{=x}%
+\sum_{\substack{n\in\mathbb{N};\\n\neq1}}\underbrace{\delta_{n,1}%
}_{\substack{=0\\\text{(since }n\neq1\text{)}}}x^{n}=x+\underbrace{\sum
_{\substack{n\in\mathbb{N};\\n\neq1}}0x^{n}}_{=0}\nonumber\\
&  =x. \label{pf.prop.fps.subs.rules.g.1}%
\end{align}


Let $g\in K\left[  \left[  x\right]  \right]  $. We must prove that $x\circ
g=g\circ x=g$. From (\ref{pf.prop.fps.subs.rules.g.1}), we obtain
$x=\sum_{n\in\mathbb{N}}\delta_{n,1}x^{n}$. Hence, Definition
\ref{def.fps.subs} (or Definition \ref{def.pol.subs}) yields $x\left[
g\right]  =\sum_{n\in\mathbb{N}}\delta_{n,1}g^{n}$ (since $\delta_{n,1}\in K$
for each $n\in\mathbb{N}$). Thus,%
\[
x\left[  g\right]  =\sum_{n\in\mathbb{N}}\delta_{n,1}g^{n}=\underbrace{\delta
_{1,1}}_{\substack{=1\\\text{(since }1=1\text{)}}}\underbrace{g^{1}}_{=g}%
+\sum_{\substack{n\in\mathbb{N};\\n\neq1}}\underbrace{\delta_{n,1}%
}_{\substack{=0\\\text{(since }n\neq1\text{)}}}g^{n}=g+\underbrace{\sum
_{\substack{n\in\mathbb{N};\\n\neq1}}0g^{n}}_{=0}=g.
\]
In other words, $x\circ g=g$ (since $x\circ g$ is a synonym for $x\left[
g\right]  $).

Next, let us write $g$ in the form $g=\sum_{n\in\mathbb{N}}g_{n}x^{n}$ for
some $g_{0},g_{1},g_{2},\ldots\in K$. Then, Definition \ref{def.fps.subs}
yields $g\left[  x\right]  =\sum_{n\in\mathbb{N}}g_{n}x^{n}=g$. Thus, $g\circ
x=g\left[  x\right]  =g$. Combining this with $x\circ g=g$, we obtain $x\circ
g=g\circ x=g$. This proves Proposition \ref{prop.fps.subs.rules} \textbf{(g)}.

\bigskip

\textbf{(b)} This appears in \cite[Theorem 7.62]{Loehr-BC} and
\cite[Proposition 2.2.2]{Brewer14}. Here is the proof:

Let $f_{1},f_{2},g\in K\left[  \left[  x\right]  \right]  $ satisfy $\left[
x^{0}\right]  g=0$. Write the FPSs $f_{1}$ and $f_{2}$ as
\begin{equation}
f_{1}=\sum_{n\in\mathbb{N}}f_{1,n}x^{n}\ \ \ \ \ \ \ \ \ \ \text{and}%
\ \ \ \ \ \ \ \ \ \ f_{2}=\sum_{n\in\mathbb{N}}f_{2,n}x^{n}
\label{pf.prop.fps.subs.rules.b.1}%
\end{equation}
with $f_{1,0},f_{1,1},f_{1,2},\ldots\in K$ and $f_{2,0},f_{2,1},f_{2,2}%
,\ldots\in K$. Thus, Definition \ref{def.fps.subs} yields%
\[
f_{1}\left[  x\right]  =\sum_{n\in\mathbb{N}}f_{1,n}x^{n}=f_{1}%
\ \ \ \ \ \ \ \ \ \ \text{and}\ \ \ \ \ \ \ \ \ \ f_{2}\left[  x\right]
=\sum_{n\in\mathbb{N}}f_{2,n}x^{n}=f_{2}%
\]
and
\[
f_{1}\left[  g\right]  =\sum_{n\in\mathbb{N}}f_{1,n}g^{n}%
\ \ \ \ \ \ \ \ \ \ \text{and}\ \ \ \ \ \ \ \ \ \ f_{2}\left[  g\right]
=\sum_{n\in\mathbb{N}}f_{2,n}g^{n}.
\]
Hence,%
\begin{align*}
f_{1}\left[  g\right]   &  =\sum_{n\in\mathbb{N}}f_{1,n}g^{n}=\sum
_{i\in\mathbb{N}}f_{1,i}g^{i}\ \ \ \ \ \ \ \ \ \ \text{and}\\
f_{2}\left[  g\right]   &  =\sum_{n\in\mathbb{N}}f_{2,n}g^{n}=\sum
_{j\in\mathbb{N}}f_{2,j}g^{j}.
\end{align*}
Multiplying these two equalities together, we find%
\begin{align}
f_{1}\left[  g\right]  \cdot f_{2}\left[  g\right]   &  =\left(  \sum
_{i\in\mathbb{N}}f_{1,i}g^{i}\right)  \left(  \sum_{j\in\mathbb{N}}%
f_{2,j}g^{j}\right)  =\sum_{i\in\mathbb{N}}\ \ \sum_{j\in\mathbb{N}%
}\underbrace{f_{1,i}g^{i}f_{2,j}g^{j}}_{\substack{=f_{1,i}f_{2,j}g^{i}%
g^{j}\\=f_{1,i}f_{2,j}g^{i+j}}}\nonumber\\
&  =\underbrace{\sum_{i\in\mathbb{N}}\ \ \sum_{j\in\mathbb{N}}}_{=\sum
_{\left(  i,j\right)  \in\mathbb{N}^{2}}}f_{1,i}f_{2,j}g^{i+j}%
=\underbrace{\sum_{\left(  i,j\right)  \in\mathbb{N}^{2}}}_{=\sum
_{n\in\mathbb{N}}\ \ \sum_{\substack{\left(  i,j\right)  \in\mathbb{N}%
^{2};\\i+j=n}}}f_{1,i}f_{2,j}g^{i+j}\nonumber\\
&  =\sum_{n\in\mathbb{N}}\ \ \sum_{\substack{\left(  i,j\right)  \in
\mathbb{N}^{2};\\i+j=n}}f_{1,i}f_{2,j}\underbrace{g^{i+j}}_{\substack{=g^{n}%
\\\text{(since }i+j=n\text{)}}}=\sum_{n\in\mathbb{N}}\ \ \sum
_{\substack{\left(  i,j\right)  \in\mathbb{N}^{2};\\i+j=n}}f_{1,i}f_{2,j}%
g^{n}\nonumber\\
&  =\sum_{n\in\mathbb{N}}\left(  \sum_{\substack{\left(  i,j\right)
\in\mathbb{N}^{2};\\i+j=n}}f_{1,i}f_{2,j}\right)  g^{n}.
\label{pf.prop.fps.subs.rules.b.3}%
\end{align}


However, we can apply the same computations to $x$ instead of $g$ (since $x$
is also an FPS with $\left[  x^{0}\right]  x=0$). Thus, we obtain%
\[
f_{1}\left[  x\right]  \cdot f_{2}\left[  x\right]  =\sum_{n\in\mathbb{N}%
}\left(  \sum_{\substack{\left(  i,j\right)  \in\mathbb{N}^{2};\\i+j=n}%
}f_{1,i}f_{2,j}\right)  x^{n}.
\]
In view of $f_{1}\left[  x\right]  =f_{1}$ and $f_{2}\left[  x\right]  =f_{2}%
$, this rewrites as%
\[
f_{1}\cdot f_{2}=\sum_{n\in\mathbb{N}}\left(  \sum_{\substack{\left(
i,j\right)  \in\mathbb{N}^{2};\\i+j=n}}f_{1,i}f_{2,j}\right)  x^{n}.
\]
Hence, Definition \ref{def.fps.subs} yields%
\[
\left(  f_{1}\cdot f_{2}\right)  \left[  g\right]  =\sum_{n\in\mathbb{N}%
}\left(  \sum_{\substack{\left(  i,j\right)  \in\mathbb{N}^{2};\\i+j=n}%
}f_{1,i}f_{2,j}\right)  g^{n}%
\]
(since $\sum_{\substack{\left(  i,j\right)  \in\mathbb{N}^{2};\\i+j=n}%
}f_{1,i}f_{2,j}\in K$ for each $n\in\mathbb{N}$). Comparing this with
(\ref{pf.prop.fps.subs.rules.b.3}), we obtain%
\[
\left(  f_{1}\cdot f_{2}\right)  \left[  g\right]  =f_{1}\left[  g\right]
\cdot f_{2}\left[  g\right]  .
\]
In other words, $\left(  f_{1}\cdot f_{2}\right)  \circ g=\left(  f_{1}\circ
g\right)  \cdot\left(  f_{2}\circ g\right)  $ (since the notation $f\circ g$
is synonymous to $f\left[  g\right]  $). This proves Proposition
\ref{prop.fps.subs.rules} \textbf{(b)}. \smallskip

\begin{fineprint}
Did you notice it? I have cheated. The above proof of Proposition
\ref{prop.fps.subs.rules} \textbf{(b)} relied on some manipulations of
infinite sums that need to be justified. Namely, we replaced \textquotedblleft%
$\sum_{i\in\mathbb{N}}\ \ \sum_{j\in\mathbb{N}}$\textquotedblright\ by
\textquotedblleft$\sum_{\left(  i,j\right)  \in\mathbb{N}^{2}}$%
\textquotedblright. This is an application of the \textquotedblleft discrete
Fubini rule\textquotedblright, and as we said above, this rule can only be
used if we know that the family $\left(  f_{1,i}f_{2,j}g^{i+j}\right)
_{\left(  i,j\right)  \in\mathbb{N}\times\mathbb{N}}$ is summable. In other
words, we need to show the following statement:

\begin{statement}
\textit{Statement 1:} For each $m\in\mathbb{N}$, all but finitely many pairs
$\left(  i,j\right)  \in\mathbb{N}\times\mathbb{N}$ satisfy $\left[
x^{m}\right]  \left(  f_{1,i}f_{2,j}g^{i+j}\right)  =0$.
\end{statement}

We shall achieve this by proving the following statement:

\begin{statement}
\textit{Statement 2:} For any three nonnegative integers $m,i,j$ with $m<i+j$,
we have $\left[  x^{m}\right]  \left(  g^{i+j}\right)  =0$.
\end{statement}

[\textit{Proof of Statement 2:} Let $m,i,j$ be three nonnegative integers with
$m<i+j$. We must show that $\left[  x^{m}\right]  \left(  g^{i+j}\right)  =0$.

We have $\left[  x^{0}\right]  g=0$. Hence, Lemma \ref{lem.fps.g=xh} (applied
to $a=g$) shows that there exists an $h\in K\left[  \left[  x\right]  \right]
$ such that $g=xh$. Consider this $h$.

Now, from $g=xh$, we obtain $g^{i+j}=\left(  xh\right)  ^{i+j}=x^{i+j}h^{i+j}%
$. However, Lemma \ref{lem.fps.first-n-coeffs-of-xna} (applied to $i+j$ and
$h^{i+j}$ instead of $k$ and $a$) shows that the first $i+j$ coefficients of
the FPS $x^{i+j}h^{i+j}$ are $0$. In other words, the first $i+j$ coefficients
of the FPS $g^{i+j}$ are $0$ (since $g^{i+j}=x^{i+j}h^{i+j}$). But $\left[
x^{m}\right]  \left(  g^{i+j}\right)  $ is one of these first $i+j$
coefficients (since $m<i+j$). Thus, we conclude that $\left[  x^{m}\right]
\left(  g^{i+j}\right)  =0$. This proves Statement 2.]

[\textit{Proof of Statement 1:} Let $m\in\mathbb{N}$. If $\left(  i,j\right)
\in\mathbb{N}\times\mathbb{N}$ is a pair satisfying $m<i+j$, then%
\begin{align*}
\left[  x^{m}\right]  \left(  f_{1,i}f_{2,j}g^{i+j}\right)   &  =f_{1,i}%
f_{2,j}\underbrace{\left[  x^{m}\right]  \left(  g^{i+j}\right)
}_{\substack{=0\\\text{(by Statement 2)}}}\ \ \ \ \ \ \ \ \ \ \left(  \text{by
(\ref{pf.thm.fps.ring.xn(la)=})}\right) \\
&  =0.
\end{align*}
Thus, all but finitely many pairs $\left(  i,j\right)  \in\mathbb{N}%
\times\mathbb{N}$ satisfy $\left[  x^{m}\right]  \left(  f_{1,i}f_{2,j}%
g^{i+j}\right)  =0$ (because all but finitely many such pairs satisfy
$m<i+j$). This proves Statement 1.]

As explained above, Statement 1 shows that the family \newline$\left(
f_{1,i}f_{2,j}g^{i+j}\right)  _{\left(  i,j\right)  \in\mathbb{N}%
\times\mathbb{N}}$ is summable, and thus our interchange of summation signs
made above is justified. This completes our proof of Proposition
\ref{prop.fps.subs.rules} \textbf{(b)}.
\end{fineprint}

\bigskip

\textbf{(c)} This follows easily from parts \textbf{(b)} and \textbf{(f)}. In
detail: Let $f_{1},f_{2},g\in K\left[  \left[  x\right]  \right]  $ be such
that $\left[  x^{0}\right]  g=0$. Assume that $f_{2}$ is invertible. Let us
first show that $f_{2}\circ g$ is invertible.

Indeed, consider the inverse $f_{2}^{-1}$ of $f_{2}$. This inverse exists
(since $f_{2}$ is invertible) and satisfies $f_{2}^{-1}\cdot f_{2}%
=\underline{1}$. Now, Proposition \ref{prop.fps.subs.rules} \textbf{(b)}
(applied to $f_{2}^{-1}$ instead of $f_{1}$) yields $\left(  f_{2}^{-1}\cdot
f_{2}\right)  \circ g=\left(  f_{2}^{-1}\circ g\right)  \cdot\left(
f_{2}\circ g\right)  $. Hence,%
\[
\left(  f_{2}^{-1}\circ g\right)  \cdot\left(  f_{2}\circ g\right)
=\underbrace{\left(  f_{2}^{-1}\cdot f_{2}\right)  }_{=\underline{1}}%
\circ\,g=\underline{1}\circ g=\underline{1}%
\]
(by Proposition \ref{prop.fps.subs.rules} \textbf{(f)}, applied to $a=1$).
Thus, the FPS $f_{2}^{-1}\circ g$ is an inverse of $f_{2}\circ g$. Hence,
$f_{2}\circ g$ is invertible. The expression $\dfrac{f_{1}\circ g}{f_{2}\circ
g}$ is therefore well-defined.

It now remains to prove that $\dfrac{f_{1}}{f_{2}}\circ g=\dfrac{f_{1}\circ
g}{f_{2}\circ g}$. To this purpose, we argue as follows: The expression
$\dfrac{f_{1}}{f_{2}}$ is well-defined, since $f_{2}$ is invertible.
Proposition \ref{prop.fps.subs.rules} \textbf{(b)} (applied to $\dfrac{f_{1}%
}{f_{2}}$ instead of $f_{1}$) yields $\left(  \dfrac{f_{1}}{f_{2}}\cdot
f_{2}\right)  \circ g=\left(  \dfrac{f_{1}}{f_{2}}\circ g\right)  \cdot\left(
f_{2}\circ g\right)  $. In view of $\dfrac{f_{1}}{f_{2}}\cdot f_{2}=f_{1}$,
this rewrites as $f_{1}\circ g=\left(  \dfrac{f_{1}}{f_{2}}\circ g\right)
\cdot\left(  f_{2}\circ g\right)  $. We can divide both sides of this equality
by $f_{2}\circ g$ (since $f_{2}\circ g$ is invertible), and thus obtain
$\dfrac{f_{1}\circ g}{f_{2}\circ g}=\dfrac{f_{1}}{f_{2}}\circ g$. In other
words, $\dfrac{f_{1}}{f_{2}}\circ g=\dfrac{f_{1}\circ g}{f_{2}\circ g}$. Thus,
Proposition \ref{prop.fps.subs.rules} \textbf{(c)} is proven.

\bigskip

\textbf{(d)} Let $f,g\in K\left[  \left[  x\right]  \right]  $ satisfy
$\left[  x^{0}\right]  g=0$. We must prove that $f^{k}\circ g=\left(  f\circ
g\right)  ^{k}$ for each $k\in\mathbb{N}$.

We prove this by induction on $k$:

\textit{Induction base:} We have $\underbrace{f^{0}}_{=\underline{1}}%
\circ\,g=\underline{1}\circ g=\underline{1}$ (by Proposition
\ref{prop.fps.subs.rules} \textbf{(f)}, applied to $a=1$). Comparing this with
$\left(  f\circ g\right)  ^{0}=\underline{1}$, we find $f^{0}\circ g=\left(
f\circ g\right)  ^{0}$. In other words, $f^{k}\circ g=\left(  f\circ g\right)
^{k}$ holds for $k=0$.

\textit{Induction step:} Let $m\in\mathbb{N}$. Assume that $f^{k}\circ
g=\left(  f\circ g\right)  ^{k}$ holds for $k=m$. We must prove that
$f^{k}\circ g=\left(  f\circ g\right)  ^{k}$ holds for $k=m+1$.

We have assumed that $f^{k}\circ g=\left(  f\circ g\right)  ^{k}$ holds for
$k=m$. In other words, we have $f^{m}\circ g=\left(  f\circ g\right)  ^{m}$.
Now,%
\begin{align*}
\underbrace{f^{m+1}}_{=f\cdot f^{m}}\circ\,g  &  =\left(  f\cdot f^{m}\right)
\circ g=\left(  f\circ g\right)  \cdot\underbrace{\left(  f^{m}\circ g\right)
}_{=\left(  f\circ g\right)  ^{m}}\\
&  \ \ \ \ \ \ \ \ \ \ \left(  \text{by Proposition \ref{prop.fps.subs.rules}
\textbf{(b)}, applied to }f_{1}=f\text{ and }f_{2}=f^{m}\right) \\
&  =\left(  f\circ g\right)  \cdot\left(  f\circ g\right)  ^{m}=\left(  f\circ
g\right)  ^{m+1}.
\end{align*}
In other words, $f^{k}\circ g=\left(  f\circ g\right)  ^{k}$ holds for
$k=m+1$. This completes the induction step. Thus, we have proven that
$f^{k}\circ g=\left(  f\circ g\right)  ^{k}$ for each $k\in\mathbb{N}$.
Proposition \ref{prop.fps.subs.rules} \textbf{(d)} is now proven.

\bigskip

\textbf{(h)} This is just a generalization of Proposition
\ref{prop.fps.subs.rules} \textbf{(a)} to (potentially) infinite sums. The
proof follows the same method, but unfortunately requires some technical
reasoning about summability. I will give the full proof for the sake of
completeness, but be warned that it contains nothing of interest.

Let $\left(  f_{i}\right)  _{i\in I}\in K\left[  \left[  x\right]  \right]
^{I}$ be a summable family of FPSs. Let $g\in K\left[  \left[  x\right]
\right]  $ be an FPS satisfying $\left[  x^{0}\right]  g=0$.

First, we shall prove that the family $\left(  f_{i}\circ g\right)  _{i\in
I}\in K\left[  \left[  x\right]  \right]  ^{I}$ is summable. Indeed, we recall
that the family $\left(  f_{i}\right)  _{i\in I}\in K\left[  \left[  x\right]
\right]  ^{I}$ is summable. In other words,%
\[
\text{for each }n\in\mathbb{N}\text{, all but finitely many }i\in I\text{
satisfy }\left[  x^{n}\right]  f_{i}=0
\]
(by the definition of \textquotedblleft summable\textquotedblright). In other
words, for each $n\in\mathbb{N}$, there exists a finite subset $I_{n}$ of $I$
such that%
\begin{equation}
\text{all }i\in I\setminus I_{n}\text{ satisfy }\left[  x^{n}\right]  f_{i}=0.
\label{pf.prop.fps.subs.rules.h.1}%
\end{equation}
Consider this subset $I_{n}$. Thus, all the sets $I_{0},I_{1},I_{2},\ldots$
are finite subsets of $I$.

Now, let $n\in\mathbb{N}$ be arbitrary. The set $I_{0}\cup I_{1}\cup\cdots\cup
I_{n}$ is a union of $n+1$ finite subsets of $I$ (because all the sets
$I_{0},I_{1},I_{2},\ldots$ are finite subsets of $I$), and thus itself is a
finite subset of $I$. Moreover,
\begin{equation}
\text{all }i\in I\setminus\left(  I_{0}\cup I_{1}\cup\cdots\cup I_{n}\right)
\text{ satisfy }\left[  x^{n}\right]  \left(  f_{i}\circ g\right)  =0.
\label{pf.prop.fps.subs.rules.h.2}%
\end{equation}


\begin{fineprint}
[\textit{Proof of (\ref{pf.prop.fps.subs.rules.h.2}):} Let $i\in
I\setminus\left(  I_{0}\cup I_{1}\cup\cdots\cup I_{n}\right)  $. We must show
that $\left[  x^{n}\right]  \left(  f_{i}\circ g\right)  =0$.

Let $m\in\left\{  0,1,\ldots,n\right\}  $. Then, $I_{m}\subseteq I_{0}\cup
I_{1}\cup\cdots\cup I_{n}$, so that $I_{0}\cup I_{1}\cup\cdots\cup
I_{n}\supseteq I_{m}$ and thus $I\setminus\underbrace{\left(  I_{0}\cup
I_{1}\cup\cdots\cup I_{n}\right)  }_{\supseteq I_{m}}\subseteq I\setminus
I_{m}$. Hence, $i\in I\setminus\left(  I_{0}\cup I_{1}\cup\cdots\cup
I_{n}\right)  \subseteq I\setminus I_{m}$. Therefore,
(\ref{pf.prop.fps.subs.rules.h.1}) (applied to $m$ instead of $n$) yields
$\left[  x^{m}\right]  f_{i}=0$.

Forget that we fixed $m$. We thus have shown that $\left[  x^{m}\right]
f_{i}=0$ for each $m\in\left\{  0,1,\ldots,n\right\}  $. In other words, the
first $n+1$ coefficients of $f_{i}$ are $0$. Hence, Lemma
\ref{lem.fps.fg-coeffs-0} (applied to $f_{i}$ and $n+1$ instead of $f$ and
$k$) shows that the first $n+1$ coefficients of $f_{i}\circ g$ are $0$.
However, $\left[  x^{n}\right]  \left(  f_{i}\circ g\right)  $ is one of these
first $n+1$ coefficients (indeed, it is the last of them); thus, this
coefficient $\left[  x^{n}\right]  \left(  f_{i}\circ g\right)  $ must be $0$.
This proves (\ref{pf.prop.fps.subs.rules.h.2}).] \smallskip
\end{fineprint}

Now, recall that $I_{0}\cup I_{1}\cup\cdots\cup I_{n}$ is a finite subset of
$I$. Hence, thanks to (\ref{pf.prop.fps.subs.rules.h.2}), we know that there
exists a finite subset $J$ of $I$ such that all $i\in I\setminus J$ satisfy
$\left[  x^{n}\right]  \left(  f_{i}\circ g\right)  =0$ (namely, $J=I_{0}\cup
I_{1}\cup\cdots\cup I_{n}$). In other words, all but finitely many $i\in I$
satisfy $\left[  x^{n}\right]  \left(  f_{i}\circ g\right)  =0$.

Forget that we fixed $n$. We thus have shown that%
\[
\text{for each }n\in\mathbb{N}\text{, all but finitely many }i\in I\text{
satisfy }\left[  x^{n}\right]  \left(  f_{i}\circ g\right)  =0.
\]
In other words, the family $\left(  f_{i}\circ g\right)  _{i\in I}\in K\left[
\left[  x\right]  \right]  ^{I}$ is summable (by the definition of
\textquotedblleft summable\textquotedblright).

It now remains to prove that $\left(  \sum_{i\in I}f_{i}\right)  \circ
g=\sum_{i\in I}f_{i}\circ g$.

For each $i\in I$, we write the FPS $f_{i}$ in the form $f_{i}=\sum
_{n\in\mathbb{N}}f_{i,n}x^{n}$ with \newline$f_{i,0},f_{i,1},f_{i,2},\ldots\in
K$. First, we shall show that
\begin{equation}
\text{the family }\left(  f_{i,m}g^{m}\right)  _{\left(  i,m\right)  \in
I\times\mathbb{N}}\text{ is summable.} \label{pf.prop.fps.subs.rules.h.3}%
\end{equation}


\begin{fineprint}
[\textit{Proof of (\ref{pf.prop.fps.subs.rules.h.3}):} Fix an $n\in\mathbb{N}%
$. Let $J$ denote the set $I_{0}\cup I_{1}\cup\cdots\cup I_{n}$. Hence, $J$ is
a union of $n+1$ finite subsets of $I$ (because all the sets $I_{0}%
,I_{1},I_{2},\ldots$ are finite subsets of $I$), and thus itself is a finite
subset of $I$. The set $J\times\left\{  0,1,\ldots,n\right\}  $ must be finite
(since it is the product of the two finite sets $J$ and $\left\{
0,1,\ldots,n\right\}  $).

Now, let $\left(  i,m\right)  \in\left(  I\times\mathbb{N}\right)
\setminus\left(  J\times\left\{  0,1,\ldots,n\right\}  \right)  $. We shall
prove that $\left[  x^{n}\right]  \left(  f_{i,m}g^{m}\right)  =0$.

We have $\left(  i,m\right)  \notin J\times\left\{  0,1,\ldots,n\right\}  $
(since $\left(  i,m\right)  \in\left(  I\times\mathbb{N}\right)
\setminus\left(  J\times\left\{  0,1,\ldots,n\right\}  \right)  $).

We note that $f_{i,m}\in K$ and thus $\left[  x^{n}\right]  \left(
f_{i,m}g^{m}\right)  =f_{i,m}\cdot\left[  x^{n}\right]  \left(  g^{m}\right)
$ (by (\ref{pf.thm.fps.ring.xn(la)=})). However, we have $\left[
x^{0}\right]  g=0$. Thus, Proposition \ref{prop.fps.subs.wd} \textbf{(a)}
(applied to $f_{i}$ and $f_{i,j}$ and $m$ instead of $f$ and $f_{j}$ and $n$)
yields that the first $m$ coefficients of the FPS $g^{m}$ are $0$. In other
words, we have
\begin{equation}
\left[  x^{k}\right]  \left(  g^{m}\right)  =0\ \ \ \ \ \ \ \ \ \ \text{for
each }k\in\left\{  0,1,\ldots,m-1\right\}  .
\label{pf.prop.fps.subs.rules.h.3.pf.5}%
\end{equation}


Now, if we have $n\in\left\{  0,1,\ldots,m-1\right\}  $, then we have $\left[
x^{n}\right]  \left(  g^{m}\right)  =0$ (by
(\ref{pf.prop.fps.subs.rules.h.3.pf.5}), applied to $k=n$) and therefore
$\left[  x^{n}\right]  \left(  f_{i,m}g^{m}\right)  =f_{i,m}\cdot
\underbrace{\left[  x^{n}\right]  \left(  g^{m}\right)  }_{=0}=0$. Hence,
$\left[  x^{n}\right]  \left(  f_{i,m}g^{m}\right)  =0$ is proved in the case
when $n\in\left\{  0,1,\ldots,m-1\right\}  $. Thus, for the rest of this proof
of $\left[  x^{n}\right]  \left(  f_{i,m}g^{m}\right)  =0$, we WLOG assume
that $n\notin\left\{  0,1,\ldots,m-1\right\}  $. Hence, $n>m-1$ (since
$n\in\mathbb{N}$), so that $n\geq m$ (since $n$ and $m$ are integers).
Therefore, $m\leq n$, so that $m\in\left\{  0,1,\ldots,n\right\}  $ and
therefore $I_{m}\subseteq I_{0}\cup I_{1}\cup\cdots\cup I_{n}=J$ (since we
defined $J$ to be $I_{0}\cup I_{1}\cup\cdots\cup I_{n}$).

If we had $i\in I_{m}$, then we would have $\left(  i,m\right)  \in
J\times\left\{  0,1,\ldots,n\right\}  $ (since $i\in I_{m}\subseteq J$ and
$m\in\left\{  0,1,\ldots,n\right\}  $), which would contradict the fact that
$\left(  i,m\right)  \notin J\times\left\{  0,1,\ldots,n\right\}  $. Thus, we
cannot have $i\in I_{m}$. Hence, $i\notin I_{m}$, so that $i\in I\setminus
I_{m}$ (since $i\in I$). Thus, (\ref{pf.prop.fps.subs.rules.h.1}) (applied to
$m$ instead of $n$) yields $\left[  x^{m}\right]  f_{i}=0$. However, from
$f_{i}=\sum_{n\in\mathbb{N}}f_{i,n}x^{n}$, we see that $\left[  x^{m}\right]
f_{i}=f_{i,m}$. Thus, $f_{i,m}=\left[  x^{m}\right]  f_{i}=0$. Consequently,
$\left[  x^{n}\right]  \left(  f_{i,m}g^{m}\right)  =\underbrace{f_{i,m}}%
_{=0}\cdot\left[  x^{n}\right]  \left(  g^{m}\right)  =0$.

Forget that we fixed $\left(  i,m\right)  $. We thus have shown that
\[
\text{all }\left(  i,m\right)  \in\left(  I\times\mathbb{N}\right)
\setminus\left(  J\times\left\{  0,1,\ldots,n\right\}  \right)  \text{ satisfy
}\left[  x^{n}\right]  \left(  f_{i,m}g^{m}\right)  =0.
\]
Therefore, all but finitely many $\left(  i,m\right)  \in I\times\mathbb{N}$
satisfy $\left[  x^{n}\right]  \left(  f_{i,m}g^{m}\right)  =0$ (since
$J\times\left\{  0,1,\ldots,n\right\}  $ is a finite subset of $I\times
\mathbb{N}$).

Forget that we fixed $n$. We thus have shown that%
\[
\text{for each }n\in\mathbb{N}\text{, all but finitely many }\left(
i,m\right)  \in I\times\mathbb{N}\text{ satisfy }\left[  x^{n}\right]  \left(
f_{i,m}g^{m}\right)  =0.
\]
In other words, the family $\left(  f_{i,m}g^{m}\right)  _{\left(  i,m\right)
\in I\times\mathbb{N}}$ is summable. This proves
(\ref{pf.prop.fps.subs.rules.h.3}).]
\end{fineprint}

Now, we have shown that the family $\left(  f_{i,m}g^{m}\right)  _{\left(
i,m\right)  \in I\times\mathbb{N}}$ is summable. Renaming the index $\left(
i,m\right)  $ as $\left(  i,n\right)  $, we thus conclude that the family
$\left(  f_{i,n}g^{n}\right)  _{\left(  i,n\right)  \in I\times\mathbb{N}}$ is
summable. The same argument (but with $g$ replaced by $x$) shows that the
family $\left(  f_{i,n}x^{n}\right)  _{\left(  i,n\right)  \in I\times
\mathbb{N}}$ is summable (since the FPS $x$ satisfies $\left[  x^{0}\right]
x=0$).

The proof of $\left(  \sum_{i\in I}f_{i}\right)  \circ g=\sum_{i\in I}%
f_{i}\circ g$ is now just a matter of computation: hen, summing the equalities
$f_{i}=\sum_{n\in\mathbb{N}}f_{i,n}x^{n}$ over all $i\in I$, we obtain%
\[
\sum_{i\in I}f_{i}=\sum_{i\in I}\ \ \sum_{n\in\mathbb{N}}f_{i,n}x^{n}%
=\sum_{n\in\mathbb{N}}\ \ \sum_{i\in I}f_{i,n}x^{n}%
\]
(here, we have been able to interchange the summation signs, since the family
$\left(  f_{i,n}x^{n}\right)  _{\left(  i,n\right)  \in I\times\mathbb{N}}$ is
summable). Thus,%
\[
\sum_{i\in I}f_{i}=\sum_{n\in\mathbb{N}}\ \ \sum_{i\in I}f_{i,n}x^{n}%
=\sum_{n\in\mathbb{N}}\left(  \sum_{i\in I}f_{i,n}\right)  x^{n}.
\]
Hence, Definition \ref{def.fps.subs} (applied to $f=\sum_{i\in I}f_{i}$)
yields%
\begin{equation}
\left(  \sum_{i\in I}f_{i}\right)  \left[  g\right]  =\sum_{n\in\mathbb{N}%
}\left(  \sum_{i\in I}f_{i,n}\right)  g^{n} \label{pf.prop.fps.subs.rules.h.5}%
\end{equation}
(since $\sum_{i\in I}f_{i,n}\in K$ for each $n\in\mathbb{N}$).

On the other hand, for each $i\in I$, we have $f_{i}\left[  g\right]
=\sum_{n\in\mathbb{N}}f_{i,n}g^{n}$ (by Definition \ref{def.fps.subs}, since
$f_{i}=\sum_{n\in\mathbb{N}}f_{i,n}x^{n}$ with $f_{i,0},f_{i,1},f_{i,2}%
,\ldots\in K$). Summing these equalities over all $i\in I$, we find%
\[
\sum_{i\in I}f_{i}\left[  g\right]  =\sum_{i\in I}\ \ \sum_{n\in\mathbb{N}%
}f_{i,n}g^{n}=\sum_{n\in\mathbb{N}}\ \ \sum_{i\in I}f_{i,n}g^{n}%
\]
(again, we have been able to interchange the summation signs, since the family
$\left(  f_{i,n}g^{n}\right)  _{\left(  i,n\right)  \in I\times\mathbb{N}}$ is
summable). Thus,%
\[
\sum_{i\in I}f_{i}\left[  g\right]  =\sum_{n\in\mathbb{N}}\ \ \sum_{i\in
I}f_{i,n}g^{n}=\sum_{n\in\mathbb{N}}\left(  \sum_{i\in I}f_{i,n}\right)
g^{n}.
\]
Comparing this with (\ref{pf.prop.fps.subs.rules.h.5}), we find $\left(
\sum_{i\in I}f_{i}\right)  \left[  g\right]  =\sum_{i\in I}f_{i}\left[
g\right]  $. In other words, $\left(  \sum_{i\in I}f_{i}\right)  \circ
g=\sum_{i\in I}f_{i}\circ g$ (since the notation $f\circ g$ is synonymous to
$f\left[  g\right]  $). Thus, Proposition \ref{prop.fps.subs.rules}
\textbf{(h)} is proven.

\bigskip

\textbf{(e)} This is \cite[Theorem 7.63]{Loehr-BC} and \cite[Proposition
2.2.5]{Brewer14}.\footnote{See also \cite[Proposition 7.6.14]{19s} for a
similar property for polynomials.} Again, let us give the proof:

Write the FPS $g$ in the form $g=\sum_{n\in\mathbb{N}}g_{n}x^{n}$ for some
$g_{0},g_{1},g_{2},\ldots\in K$. Then, $g_{0}=\left[  x^{0}\right]  g=0$.
Moreover, $g\circ h=g\left[  h\right]  =\sum_{n\in\mathbb{N}}g_{n}h^{n}$ (by
Definition \ref{def.fps.subs}, because $g=\sum_{n\in\mathbb{N}}g_{n}x^{n}$
with $g_{0},g_{1},g_{2},\ldots\in K$). But Proposition \ref{prop.fps.subs.wd}
\textbf{(c)} (applied to $g$, $h$ and $g_{n}$ instead of $f$, $g$ and $f_{n}$)
yields $\left[  x^{0}\right]  \left(  \sum_{n\in\mathbb{N}}g_{n}h^{n}\right)
=g_{0}=0$. In view of $g\circ h=\sum_{n\in\mathbb{N}}g_{n}h^{n}$, this
rewrites as $\left[  x^{0}\right]  \left(  g\circ h\right)  =0$. Hence, the
composition $f\circ\left(  g\circ h\right)  $ is well-defined.

It remains to show that $\left(  f\circ g\right)  \circ h=f\circ\left(  g\circ
h\right)  $.

Write the FPS $f$ in the form $f=\sum_{n\in\mathbb{N}}f_{n}x^{n}$ for some
$f_{0},f_{1},f_{2},\ldots\in K$. Thus, Definition \ref{def.fps.subs} yields%
\[
f\left[  g\right]  =\sum_{n\in\mathbb{N}}f_{n}g^{n}%
\ \ \ \ \ \ \ \ \ \ \text{and}\ \ \ \ \ \ \ \ \ \ f\left[  g\circ h\right]
=\sum_{n\in\mathbb{N}}f_{n}\cdot\left(  g\circ h\right)  ^{n}.
\]
Moreover, the family $\left(  f_{n}g^{n}\right)  _{n\in\mathbb{N}}$ is
summable (by Proposition \ref{prop.fps.subs.wd} \textbf{(b)}). Hence,
Proposition \ref{prop.fps.subs.rules} \textbf{(h)} (applied to $\left(
f_{n}g^{n}\right)  _{n\in\mathbb{N}}$ and $h$ instead of $\left(
f_{i}\right)  _{i\in I}$ and $g$) yields that the family $\left(  \left(
f_{n}g^{n}\right)  \circ h\right)  _{n\in\mathbb{N}}\in K\left[  \left[
x\right]  \right]  ^{\mathbb{N}}$ is summable as well and that we have
\begin{equation}
\left(  \sum_{n\in\mathbb{N}}f_{n}g^{n}\right)  \circ h=\sum_{n\in\mathbb{N}%
}\left(  f_{n}g^{n}\right)  \circ h. \label{pf.prop.fps.subs.rules.e.3}%
\end{equation}
In view of $f\circ g=f\left[  g\right]  =\sum_{n\in\mathbb{N}}f_{n}g^{n}$, we
can rewrite (\ref{pf.prop.fps.subs.rules.e.3}) as%
\begin{equation}
\left(  f\circ g\right)  \circ h=\sum_{n\in\mathbb{N}}\left(  f_{n}%
g^{n}\right)  \circ h. \label{pf.prop.fps.subs.rules.e.4}%
\end{equation}
However, for each $n\in\mathbb{N}$, we have $f_{n}g^{n}=\underline{f_{n}}\cdot
g^{n}$ (by Theorem \ref{thm.fps.ring} \textbf{(d)}, applied to $\lambda=f_{n}$
and $\mathbf{a}=g^{n}$) and thus%
\begin{align}
\left(  f_{n}g^{n}\right)  \circ h  &  =\left(  \underline{f_{n}}\cdot
g^{n}\right)  \circ h=\underbrace{\left(  \underline{f_{n}}\circ h\right)
}_{\substack{=\underline{f_{n}}\\\text{(by Proposition
\ref{prop.fps.subs.rules} \textbf{(f)},}\\\text{applied to }f_{n}\text{ and
}h\\\text{instead of }a\text{ and }g\text{)}}}\cdot\underbrace{\left(
g^{n}\circ h\right)  }_{\substack{=\left(  g\circ h\right)  ^{n}\\\text{(by
Proposition \ref{prop.fps.subs.rules} \textbf{(d)},}\\\text{applied to
}g\text{ and }h\\\text{instead of }f\text{ and }g\text{)}}}\nonumber\\
&  \ \ \ \ \ \ \ \ \ \ \ \ \ \ \ \ \ \ \ \ \left(
\begin{array}
[c]{c}%
\text{by Proposition \ref{prop.fps.subs.rules} \textbf{(b)},}\\
\text{applied to }\underline{f_{n}}\text{, }g^{n}\text{ and }h\text{ instead
of }f_{1}\text{, }f_{2}\text{ and }g
\end{array}
\right) \nonumber\\
&  =\underline{f_{n}}\cdot\left(  g\circ h\right)  ^{n}=f_{n}\cdot\left(
g\circ h\right)  ^{n} \label{pf.prop.fps.subs.rules.e.5}%
\end{align}
(since Theorem \ref{thm.fps.ring} \textbf{(d)} (applied to $\lambda=f_{n}$ and
$\mathbf{a}=\left(  g\circ h\right)  ^{n}$) yields $f_{n}\cdot\left(  g\circ
h\right)  ^{n}=\underline{f_{n}}\cdot\left(  g\circ h\right)  ^{n}$). Thus,
(\ref{pf.prop.fps.subs.rules.e.4}) becomes%
\begin{align*}
\left(  f\circ g\right)  \circ h  &  =\sum_{n\in\mathbb{N}}\underbrace{\left(
f_{n}g^{n}\right)  \circ h}_{\substack{=f_{n}\cdot\left(  g\circ h\right)
^{n}\\\text{(by (\ref{pf.prop.fps.subs.rules.e.5}))}}}=\sum_{n\in\mathbb{N}%
}f_{n}\cdot\left(  g\circ h\right)  ^{n}\\
&  =f\left[  g\circ h\right]  \ \ \ \ \ \ \ \ \ \ \left(  \text{since
}f\left[  g\circ h\right]  =\sum_{n\in\mathbb{N}}f_{n}\cdot\left(  g\circ
h\right)  ^{n}\right) \\
&  =f\circ\left(  g\circ h\right)  .
\end{align*}
This completes the proof of Proposition \ref{prop.fps.subs.rules} \textbf{(e)}.
\end{proof}

\begin{example}
Let us use Proposition \ref{prop.fps.subs.rules} \textbf{(c)} to justify the
equality (\ref{eq.exa.fps.subs.fibonacci.subs-plausible}) that we used in
Example \ref{exa.fps.subs.fibonacci}. Indeed, we know that the FPS $1-x$ is
invertible. Thus, applying Proposition \ref{prop.fps.subs.rules} \textbf{(c)}
to $f_{1}=1$ and $f_{2}=1-x$ and $g=x+x^{2}$, we obtain%
\[
\dfrac{1}{1-x}\circ\left(  x+x^{2}\right)  =\dfrac{1\circ\left(
x+x^{2}\right)  }{\left(  1-x\right)  \circ\left(  x+x^{2}\right)  }.
\]
Using the notation $f\left[  g\right]  $ instead of $f\circ g$, we can rewrite
this as%
\[
\dfrac{1}{1-x}\left[  x+x^{2}\right]  =\dfrac{1\left[  x+x^{2}\right]
}{\left(  1-x\right)  \left[  x+x^{2}\right]  }.
\]
In view of $1\left[  x+x^{2}\right]  =1$ and $\left(  1-x\right)  \left[
x+x^{2}\right]  =1-\left(  x+x^{2}\right)  =1-x-x^{2}$, this rewrites as
$\dfrac{1}{1-x}\left[  x+x^{2}\right]  =\dfrac{1}{1-x-x^{2}}$. Thus,
(\ref{eq.exa.fps.subs.fibonacci.subs-plausible}) is proved.
\end{example}

Let us summarize: If $f,g\in K\left[  \left[  x\right]  \right]  $ are two
FPSs, then the composition $f\left[  g\right]  $ is not always defined.
However, it is defined at least in the following two cases:

\begin{itemize}
\item in the case when $f$ is a polynomial (that is, $f\in K\left[  x\right]
$), and

\item in the case when $g$ has constant term $0$ (that is, $\left[
x^{0}\right]  g=0$).
\end{itemize}

This justifies some more of the things we did back in Section
\ref{sec.gf.exas}; in particular, Example 1 from that section is now fully
justified. But we still have not defined (e.g.) the square root of an FPS,
which we used in Example 2.

Before I explain square roots, let me quickly survey differentiation of FPSs.
